\chapter{The Small-Signal Regime: Floquet Scattering}
\label{chap:floquet}

\section{The regime}

\Cref{chap:pumphb} produced $\Xvec$, the periodic steady state of the pumped
circuit. Now inject a signal so weak that it does not perturb that state. The
question this chapter answers is: what linear system does the signal see?

The answer is the defining structural insight of the whole solver. With the pump
fixed, the signal sees a circuit whose inductances vary \emph{periodically in
time} but which is otherwise \emph{linear}. That is a linear time-periodic (LTP)
system, and Floquet theory applies exactly.

\begin{definition}[small-signal regime]
The regime in which the signal amplitude is small enough that (i) the pump
solution $\Xvec$ is unchanged by its presence, and (ii) the response is linear in
the signal amplitude. Both conditions fail at large signal power; that is the
subject of \cref{chap:saturation}.
\end{definition}

This is why a gain map is affordable and a compression sweep is not. In this
regime the sidebands are a \emph{separate linear system} that can be factored
once per frequency; they do not act back on the pump, so no outer nonlinear
iteration is required.

\section{Floquet structure}

\subsection{The time-periodic tangent}

Linearising \eqref{eq:governing} about the pumped state, a perturbation
$\delta\xvec$ obeys
\begin{equation}
  \Cmat\,\delta\ddot{\xvec}
  + \Gmat\,\delta\dot{\xvec}
  + \Kmat\,\delta\xvec
  + \Bphi\,\diag\bigl(\gamma(t)\bigr)\,\Bphi^{\mathsf T}\delta\xvec
  = \delta\isrc,
  \label{eq:linearised}
\end{equation}
with the junction tangent of \eqref{eq:gamma-t},
\begin{equation}
  \gamma_b(t) = \frac{I_{c,b}}{\rphi}
                \cos\!\left(\frac{\psi_{p,b}(t)}{\rphi}\right),
\end{equation}
evaluated on the \emph{converged pump waveform}. Because $\psi_p$ is
$T$-periodic, so is $\gamma$.

\subsection{Sideband coupling}

Expand $\gamma$ in a Fourier series,
\begin{equation}
  \gamma(t) = \sum_{\ell\in\mathbb{Z}} \ghat_\ell\,e^{i\ell\wp t},
  \qquad
  \ghat_\ell = \frac{1}{N_t}\sum_{r}\gamma(t_r)e^{-i\ell\wp t_r},
\end{equation}
and note that a term $\ghat_\ell e^{i\ell\wp t}$ multiplying a signal component
at $\ws$ produces a component at $\ws + \ell\wp$. Floquet's theorem states that
solutions of \eqref{eq:linearised} therefore take the form
\begin{equation}
  \delta\xvec(t) = \sum_{m\in\mathbb{Z}} \Phi_m\,e^{i(\ws + m\wp)t},
\end{equation}
a \emph{ladder} of sidebands spaced by the pump frequency.

\begin{definition}[sideband frequency]
\begin{equation}
  \omega_m = \ws + m\wp,\qquad m\in\mathbb{Z}.
\end{equation}
$m=0$ is the signal. For four-wave mixing the idler is $m=-2$, since
$2\wp = \ws + \wi$ gives $\wi = 2\wp - \ws = -(\ws - 2\wp)$; the negative
frequency is represented as the conjugate of the $m=-2$ component.
\end{definition}

\subsection{The conversion matrix}

Substituting the ladder into \eqref{eq:linearised} and projecting onto
$e^{-i\omega_m t}$ gives, for each $m$,
\begin{equation}
  \Dmat(\omega_m)\,\Phi_m
  + \sum_{q}\Khat_{m-q}\,\Phi_q
  = \delta\Svec_m ,
\end{equation}
with $\Khat_\ell$ as in \eqref{eq:khat}. Truncating the ladder to
$m\in[-S,S]$ gives a finite block matrix

\begin{equation}
  \boxed{\;
  \mat{A}
  = \blockdiag\bigl(\Dmat(\omega_{-S}),\dots,\Dmat(\omega_{S})\bigr)
  \;+\;
  \bigl[\,\Khat_{m-q}\,\bigr]_{m,q=-S}^{S},
  \;}
  \label{eq:conversion-matrix}
\end{equation}
of dimension $(2S+1)n$. The signal problem is then the \emph{single linear
solve}
\begin{equation}
  \mat{A}\,\vect{y} = \vect{b},
  \label{eq:signal-solve}
\end{equation}
with $\vect{b}$ placing the source current at the signal sideband and source
node.

\begin{implnote}[assembly]
\code{assemble\_conversion\_matrix} in \file{signal/floquet.py} builds
\eqref{eq:conversion-matrix} block by block, adding $\Dmat(\omega_m)$ on the
diagonal and $\Khat_{m-q}$ everywhere. A cached variant,
\code{assemble\_conversion\_matrix\_from\_base}, precomputes the $\Khat$
block structure once per pump point --- it does not depend on $\ws$ --- and adds
only the diagonal $\Dmat(\omega_m)$ blocks per signal frequency. That is the
dominant saving in a frequency sweep.
\end{implnote}

\subsection{The conjugate-symmetry check}

\begin{proposition}
For a genuinely real pump waveform,
\begin{equation}
  \ghat_{-\ell} = \overline{\ghat_{\ell}}
  \qquad\text{exactly.}
  \label{eq:conj-symmetry}
\end{equation}
\end{proposition}
\begin{proof}
$\gamma(t)$ is real, so its Fourier coefficients satisfy Hermitian symmetry.
\end{proof}

\begin{implnote}[\code{gamma\_hat\_summary.csv}]
This is a powerful diagnostic because it is an \emph{exact} identity, not an
approximate one. \code{write\_gamma\_hat\_summary} emits, per sideband $\ell$:
\code{ell, nbranches, l2\_abs, l2\_abs\_over\_zero\_l2, max\_abs, mean\_abs,
mean\_real, mean\_imag, conj\_symmetry\_rel\_err}, with
\begin{equation}
  \texttt{conj\_symmetry\_rel\_err}
  = \frac{\|\ghat_\ell - \overline{\ghat_{-\ell}}\|_2}{\|\ghat_\ell\|_2}.
\end{equation}
Any nonzero value means the reconstructed pump waveform is not real ---
typically a phase-convention or basis-metadata mismatch. It does \emph{not} mean
the gain is slightly off; it means the pump is wrong.
\end{implnote}

\section{Gain}

\subsection{Three definitions, and which to use}

The solver reports three distinct gain quantities. They answer different
questions and are not interchangeable.

\begin{center}
\begin{tabular}{@{}l l p{0.44\textwidth}@{}}
\toprule
Field & Definition & Meaning \\
\midrule
\code{gain\_db}
& $20\log_{10}|S_{21}|$
& Absolute transducer gain from the port S-parameter
  \eqref{eq:s-from-response}. \\
\addlinespace
\code{gain\_vs\_off\_db}
& $10\log_{10}\bigl|V_{\mathrm{on}}/V_{\mathrm{off}}\bigr|^{2}$
& Ratio to the same circuit with the pump off. This is what an experiment
  measures as ``gain''. \\
\addlinespace
\code{gain\_vs\_pumpdiag\_db}
& $10\log_{10}\bigl|V_{\mathrm{on}}/V_{\mathrm{pumpdiag}}\bigr|^{2}$
& Ratio to a reference including the pump-induced \emph{mean} stiffness
  $\ghat_0$ but no sideband conversion. Isolates conversion gain from the
  pump's static softening of the line. \\
\bottomrule
\end{tabular}
\end{center}

\begin{caveat}[do not subtract absolute S values]
\code{gain\_vs\_off\_db} must be formed as a ratio of voltage-per-input-current
responses, which is what the Floquet normalisation gives. The absolute one-port
$S$ of \eqref{eq:s-from-response} is \emph{affine} in the response --- it
carries the $-1$ reflection subtraction --- so differencing two absolute $S$
values does not give a gain ratio. This was a real defect and is now the
documented contract.
\end{caveat}

\subsection{Baselines}

Both baselines are single-block solves at the signal frequency alone, with no
sideband ladder:
\begin{equation}
  \left[\Dmat(\ws) + \mat{D}_{\mathrm{extra}}\right]\vect{y} = \vect{b},
\end{equation}
with $\mat{D}_{\mathrm{extra}} = \Khat_0^{\mathrm{off}}$ (junction inductance
only, pump off) or $\Khat_0$ (pump-averaged stiffness). They are cheap --- one
$n\times n$ solve each against the $(2S+1)n$ coupled solve.

\subsection{Idler}

If the idler sideband is retained, its output voltage is extracted at its own
signed frequency,
\begin{equation}
  V_{\mathrm{idler}} = i(\ws + m_i\wp)\,\Phi_{m_i}\big|_{\text{out node}},
\end{equation}
and reported relative to the pump-off signal. Using $\ws$ rather than the idler's
own frequency here would misscale the idler voltage.

\section{Solution quality}

\begin{implnote}[status is earned, not assumed]
After solving \eqref{eq:signal-solve} the solver forms the explicit linear
residual
\begin{lstlisting}
r = A @ y - b
linear_rel_residual = norm(r) / max(norm(b), 1e-300)
\end{lstlisting}
and downgrades the status from \code{VALID\_SOLVED} to \code{CHECK} if the
relative residual exceeds $10^{-7}$ or the gain is non-finite. A direct solver
returning without an exception is not by itself evidence that the system was
well conditioned.
\end{implnote}

\section{Backends}

\paragraph{Direct.} Factor the full $(2S+1)n$ system \eqref{eq:conversion-matrix}.

\paragraph{Schur.} Eliminate interior nodes first (\cref{sec:schur}) and solve
the reduced coupled system. Faster per solve, but the Schur complement is denser
than the original block, and at map scale that density can exhaust memory.

\begin{caveat}[choose \code{direct} for large campaigns]
The Schur signal backend densifies. For a $50\times50$ map with many signal
frequencies this has caused out-of-memory failures; \code{direct} is the safe
default for campaigns, with \code{schur} reserved for single-point work.
\end{caveat}

\section{Passive response}

With the pump off, \eqref{eq:conversion-matrix} collapses to the single block
\begin{equation}
  \Dmat(\omega) + \Bphi\,\diag\!\left(\frac{\Ic}{\rphi}\right)\Bphi^{\mathsf T},
\end{equation}
since $\ghat_0 = \Ic/\rphi$ at zero flux and all other $\ghat_\ell$ vanish.

\begin{implnote}[\code{passive\_s\_matrix}]
This is a genuine linearised S-parameter, not a lumped approximation: the
Josephson inductance is carried through $\ghat_0$ exactly as in the pumped case.
The convention is
\begin{equation}
  S[\text{frequency},\ \text{output port},\ \text{source port}],
\end{equation}
and the four directional traces plotted by the passive workflow ($S_{11}$,
$S_{21}$, $S_{31}$, $S_{41}$) are all excited from port 1.
\end{implnote}

The passive response is the first sanity check on any new design: it must be
finite, continuous, and passive ($|S_{21}|\le1$) across the band before any
pumped result is meaningful.

\section{Quantum efficiency}

\subsection{The ladder-operator basis}

Quantum efficiency asks what fraction of the power arriving in an output mode
came from the intended input mode rather than from vacuum or from another mode.
For a scattering matrix in the \emph{photon ladder-operator} basis,
\begin{equation}
  \eta_{ij} = \frac{|S_{ij}|^{2}}
                   {\sum_k \left(|S_{ik}|^{2} + |S^{\mathrm{noise}}_{ik}|^{2}\right)},
  \label{eq:qe}
\end{equation}
and the ideal (best achievable) efficiency for a given $|S|$ is
\begin{equation}
  \eta^{\mathrm{ideal}}
  = \begin{cases}
      1, & |S|\le 1,\\[3pt]
      \dfrac{1}{2 - 1/|S|^{2}}, & |S| > 1 .
    \end{cases}
  \label{eq:qe-ideal}
\end{equation}

\begin{caveat}[basis conversion is mandatory]
Equations \eqref{eq:qe}--\eqref{eq:qe-ideal} require $S$ in the ladder-operator
basis. What \code{solve\_gain\_one} returns is the \emph{classical
voltage-ratio} $S$. Signal and idler sit at different frequencies, so converting
requires the Manley--Rowe reweighting
\begin{equation}
  S^{\mathrm{ladder}}_{mn}
  = S^{\mathrm{classical}}_{mn}\sqrt{\frac{\omega_n}{\omega_m}} .
  \label{eq:ladder-weights}
\end{equation}
Applying \eqref{eq:qe} to the classical matrix silently gives wrong numbers.
\end{caveat}

\subsection{The unitarity check and its truncation}

For a lossless non-degenerate parametric amplifier, unitarity in the ladder
basis requires
\begin{equation}
  |S_{ss}|^{2} - |S_{is}|^{2} = 1 .
  \label{eq:unitarity}
\end{equation}

\begin{measured}[$2\times2$ truncation]
Building the $2\times2$ $[\text{signal},\text{idler}]$ sub-matrix from two
\code{solve\_gain\_one\_schur} calls and testing \eqref{eq:unitarity}:
\begin{center}
\begin{tabular}{@{}l r@{}}
\toprule
Device & Deviation from unitarity \\
\midrule
\code{ipm\_2c\_fixed} & $\approx\SI{2}{\percent}$ \\
\code{jc\_fqjtwpa} & $\approx\SI{0.4}{\percent}$ \\
\code{jc\_jtwpa} & $\approx\SI{47}{\percent}$ \\
\bottomrule
\end{tabular}
\end{center}
\end{measured}

\begin{remark}[the $\SI{47}{\percent}$ is expected, not a bug]
\code{jc\_jtwpa} is solved with ten sidebands. A $2\times2$ truncation discards
the power leaking into the other eight, and \eqref{eq:unitarity} is a statement
about the \emph{full} scattering matrix. A genuine unitarity test requires the
complete multi-sideband $S$; the $2\times2$ version is a compact demonstration
of the basis conversion, nothing more. That the two devices with weaker sideband
content pass to a few percent is consistent with exactly this reading.
\end{remark}

\begin{implnote}[\code{calc\_qe} / \code{calc\_qe\_ideal}]
Direct ports of the corresponding routines in a reference Julia code,
vectorised. The Julia sources provide separate in-place (\code{!}-suffixed) and
two-pass variants written by hand for cache efficiency; NumPy's row-sum already
gives that behaviour, so only the single vectorised form is implemented. The
docstring examples of the original are ported verbatim as tests.
\end{implnote}

\section{What the small-signal stage cannot answer}

Everything in this chapter assumes the signal does not disturb the pump. That
assumption is built into the structure of \eqref{eq:conversion-matrix}: the
$\Khat_\ell$ blocks are computed from the \emph{pump-only} waveform and are held
fixed.

Consequently the small-signal stage cannot:
\begin{itemize}
\item predict compression at any signal power;
\item account for pump depletion;
\item detect a change of phase matching caused by the signal;
\item say anything about the stability of a large-signal state.
\end{itemize}

Each requires solving the pump and the signal together. That is
\cref{chap:saturation}.

\begin{remark}[the small-signal limit as a validation target]
The relationship runs the other way too. In the limit of vanishing signal power,
the multitone solution of \cref{chap:saturation} must reproduce the Floquet
result of this chapter at the same sideband count. That is a genuine internal
consistency check --- the two calculations share almost no code --- and it
remains one of the few available checks on the saturation machinery
(\cref{chap:validation}).
\end{remark}
